\section*{अध्याय \ref{ch:g8:the-atmosphere} --- वायुमंडल}

\begin{solution}{exo:g8:the-atmosphere:1}
(क) \omterm{def:g8:the-atmosphere:layers}{समतापमंडल}; (ख) \omterm{def:g8:the-atmosphere:layers}{क्षोभमंडल}; (ग) \omterm{def:g8:the-atmosphere:layers}{मध्यमंडल}; (घ) \omterm{def:g8:the-atmosphere:layers}{तापमंडल}।
\end{solution}

\begin{solution}{exo:g8:the-atmosphere:2}
(क) \qty{15}{\degreeCelsius}। (ख) लगभग \qty{-56}{\degreeCelsius}।
(ग) लगभग \qty{-3}{\degreeCelsius}। (घ) लगभग \qty{-86}{\degreeCelsius}।
(ङ) \qty{11}{km} से लगभग \qty{20}{km} तक (और \qty{50}{km} के पास तथा
\qty{86}{km} के ठीक ऊपर कुछ किलोमीटरों में)।
\end{solution}

\begin{solution}{exo:g8:the-atmosphere:3}
(क) ग़लत: वह \omterm{def:g8:the-atmosphere:layers}{क्षोभमंडल} में गिरता है, \omterm{def:g8:the-atmosphere:layers}{समतापमंडल} में बढ़ता है, \omterm{def:g8:the-atmosphere:layers}{मध्यमंडल} में गिरता है
और \omterm{def:g8:the-atmosphere:layers}{तापमंडल} में फिर बढ़ता है। (ख) सही। (ग) ग़लत: वह \omterm{def:g8:the-atmosphere:layers}{समतापमंडल} में है। (घ) सही।
\end{solution}

\begin{solution}{exo:g8:the-atmosphere:4}
(क) \omterm{def:g8:the-atmosphere:layers}{क्षोभमंडल}। (ख) \omterm{def:g8:the-atmosphere:layers}{समतापमंडल} (मानक वायुमंडल में \omterm{def:g8:the-atmosphere:layers}{क्षोभसीमा} \qty{11}{km} पर है)।
(ग) \omterm{def:g8:the-atmosphere:layers}{समतापमंडल}। (घ) \omterm{def:g8:the-atmosphere:layers}{मध्यमंडल}।
\end{solution}

\begin{solution}{exo:g8:the-atmosphere:5}
(क) $20 - 6.5 \times 2 = 20 - 13 = \qty{7}{\degreeCelsius}$। (ख) $20 -
6.5 \times 3 = 20 - 19.5 = \qty{0.5}{\degreeCelsius}$। (ग) हवा को 20 डिग्री ठंडा होना
है: $20 \div 6.5 \approx 3.1$, इसलिए लगभग
\qty{3.1}{km} पर।
\end{solution}

\begin{solution}{exo:g8:the-atmosphere:6}
(क) $3 - 0.5 = \qty{2.5}{km}$। (ख) $6.5 \times 2.5 = 16.25$: लगभग 16
डिग्री। (ग) $18 - 16.25 \approx \qty{2}{\degreeCelsius}$।
\end{solution}

\begin{solution}{exo:g8:the-atmosphere:7}
(क) \qty{500}{hPa}, \qty{250}{hPa} और \qty{125}{hPa}। (ख) ज़मीन के पास नियम
अच्छा है (\qty{500}{hPa} बनाम \qty{506}{hPa})। और ऊपर वह बहुत ज़्यादा देता है:
\qty{250}{hPa} बनाम \qty{227}{hPa}, और \qty{125}{hPa} बनाम \qty{96}{hPa}।
आधे होने का नियम केवल एक मोटा
अनुमान है।
\end{solution}

\begin{solution}{exo:g8:the-atmosphere:8}
(क) $10~\text{hPa} = 10^1$~hPa; $10^3 \div 10^1 = 10^2$: सौ गुना छोटा।
(ख) $0.01~\text{hPa} = 10^{-2}$~hPa। (ग) $10^3 \div 10^{-2} =
10^5$: एक लाख
गुना छोटा।
\end{solution}

\begin{solution}{exo:g8:the-atmosphere:9}
(क) $6 \div 1.5 = 4$ गुना चौड़ा। (ख) ऊपर उठते हुए उसके चारों ओर की हवा कम और कम
दबाती है, इसलिए अंदर की गैस फैलती है और गुब्बारा तब तक बढ़ता है जब तक उसकी झिल्ली फट
नहीं जाती। (ग) गुब्बारा फटने के बाद पैराशूट रेडियोसोंड के गिरने को धीमा करता है, ताकि
वह किसी को चोट न पहुँचाए या किसी चीज़ को नुक़सान न
करे।
\end{solution}

\begin{solution}{exo:g8:the-atmosphere:10}
(क) $20 \div 7 \approx 2.9$ गुना मोटा। (ख) \omterm{def:g3:continents-and-oceans:poles}{ध्रुवों} के ज़्यादा पास: \qty{11}{km} से
दूरी \qty{4}{km} है, \qty{7}{km} वाली ध्रुवीय \omterm{def:g8:the-atmosphere:layers}{क्षोभसीमा} तक, पर \qty{9}{km},
\qty{20}{km} वाली भूमध्यरेखीय \omterm{def:g8:the-atmosphere:layers}{क्षोभसीमा} तक (वह मध्य \omterm{def:g6:latitude-longitude:latitude}{अक्षांशों} का वर्णन करता है,
\omterm{def:g3:continents-and-oceans:equator}{भूमध्य रेखा} का नहीं)। (ग) शुरुआत में \omterm{def:g8:the-atmosphere:layers}{क्षोभसीमा} लगभग \qty{7}{km} पर है, विमान के
नीचे: विमान \omterm{def:g8:the-atmosphere:layers}{समतापमंडल} में है। अंत में \omterm{def:g8:the-atmosphere:layers}{क्षोभसीमा} लगभग \qty{20}{km} पर है: विमान
\omterm{def:g8:the-atmosphere:layers}{क्षोभमंडल} में है।
\end{solution}

\begin{solution}{exo:g8:the-atmosphere:11}
(क) $1 - 506 \div 1013 \approx 0.50$: \textbf{50~\%}। (ख) $1 - 227 \div
1013 \approx 0.78$: \textbf{78~\%}। (ग) $1 - 12 \div 1013 \approx 0.99$:
\textbf{99~\%}। (घ) लगभग 1~\%।
\end{solution}

\begin{solution}{exo:g8:the-atmosphere:12}
(क) $6371 \div 15 \approx 425$~km हर सेंटीमीटर के लिए। (ख) $100 \div 425
\approx 0.24$~cm, लगभग \qty{2.4}{mm}। (ग) $11 \div 425 \approx
0.026$~cm,
लगभग \qty{0.26}{mm}। (घ) $0.26 \div 0.1 \approx 2.6$: काग़ज़ की दो या तीन
शीटें।
\end{solution}

\begin{solution}{pb:g8:the-atmosphere:1}
\textbf{1.} $300 \times 10 = 3000$~m।

\textbf{2.} $6000 \div 300 = 20$ मिनट; $\num{30000} \div 300 = 100$
मिनट।

\textbf{3.} 100 मिनट, यानी 1~घंटा और 40~मिनट।

\textbf{4.} \qty{8}{km} पर: \omterm{def:g8:the-atmosphere:layers}{क्षोभमंडल} (तापमान अब भी गिर रहा है)। \qty{30}{km}
पर: \omterm{def:g8:the-atmosphere:layers}{समतापमंडल}।

\textbf{5.} $15 - (-50) = 65$ डिग्री।

\textbf{6.} $65 \div 10 = \qty{6.5}{\degreeCelsius}$ प्रति किलोमीटर।

\textbf{7.} \qty{12}{km} से \qty{20}{km} तक वह \qty{-56}{\degreeCelsius} पर रहता
है। ठंडा होना \qty{10}{km} और \qty{12}{km} के बीच रुकता है: \omterm{def:g8:the-atmosphere:layers}{क्षोभसीमा} लगभग
\qty{11}{km} पर है।

\textbf{8.} वह बढ़ता है: \qty{-55}{\degreeCelsius}, \qty{22}{km} पर;
\qty{-47}{\degreeCelsius}, \qty{30}{km} पर। गुब्बारा \omterm{def:g8:the-atmosphere:layers}{समतापमंडल} में है, जिसे सूर्य का
पराबैंगनी प्रकाश सोखने वाली ओज़ोन गरम करती
है।

\textbf{9.} $1013 \div 2 \approx \qty{507}{hPa}$; यह \qty{4}{km} (\qty{617}{hPa})
और \qty{6}{km} (\qty{472}{hPa}) के बीच पहुँचा जाता है।

\textbf{10.} $1013 \div 4 \approx \qty{253}{hPa}$।

\textbf{11.} \qty{10}{km} (\qty{265}{hPa}) और \qty{12}{km} (\qty{194}{hPa})
के बीच।

\textbf{12.} \qty{10}{km} से, दाब को $265 - 253 =
12$~hPa और गिरना है, उन 71~hPa
में से जो वह 2~km में खोता है: $12 \div 71 \times 2
\approx 0.34$~km। \omterm{def:g4:world-physical-map:altitude}{समुद्र तल}
के मान का एक चौथाई दाब लगभग \textbf{\qty{10.3}{km}} पर
है।
\end{solution}
